% Generated draft: requires mathematical and source review.
\documentclass[11pt]{article}
\usepackage[margin=25mm]{geometry}
\usepackage{fontspec}
\setmainfont{Latin Modern Roman}
\usepackage{amsmath,amssymb,mathtools,unicode-math}
\setmathfont{Latin Modern Math}
\usepackage{enumitem,needspace,xurl,hyperref}
\setlength{\parindent}{0pt}
\setlength{\parskip}{5pt}
\setlength{\emergencystretch}{3em}
\newcommand{\N}{\mathbb{N}}
\newcommand{\Z}{\mathbb{Z}}
\newcommand{\Q}{\mathbb{Q}}
\newcommand{\R}{\mathbb{R}}
\newcommand{\C}{\mathbb{C}}
\newcommand{\F}{\mathbb{F}}
\newcommand{\A}{\mathbb{A}}
\newcommand{\PP}{\mathbb{P}}
\newcommand{\E}{\mathbb{E}}
\newcommand{\eps}{\varepsilon}
\newcommand{\dd}{\,\mathrm{d}}
\begin{document}
\section*{UnsolvedMath ID: 30004519}
\textbf{Source problem code:} catalogue entry
\subsection{Definitions and mathematical statement}
The supplied question asks whether bistellar-flip graphs of triangulations of every point set are connected in dimensions \(3\) and \(4\). The record does not explicitly specify finiteness, affine dimension, general position, or whether every point must be used as a vertex. The following makes two possible finite, full-dimensional interpretations precise; identifying the intended convention requires source review.

Let \(d\in\{3,4\}\) and let \(A\subset\R^d\) be finite with \(\operatorname{aff}(A)=\R^d\). A geometric triangulation of \(\operatorname{conv}(A)\) is a finite geometric simplicial complex whose union is \(\operatorname{conv}(A)\). Consider two conventions: \(\mathcal T_{\mathrm{all}}(A)\) consists of such triangulations with vertex set exactly \(A\), while \(\mathcal T_{\mathrm{sub}}(A)\) permits any vertex set contained in \(A\). No general-position assumption is imposed in these candidate formulations.

Here is an explicit convention for bistellar flips, identifying geometric simplices with their vertex sets. A circuit \(C\subseteq A\) is a minimally affinely dependent set. Its signed partition \(C=C_+\sqcup C_-\) is determined, up to exchanging signs, by an affine dependence with positive coefficients on \(C_+\) and negative coefficients on \(C_-\). Its two circuit triangulations are
\[
 K_+=\{\sigma\subseteq C:C_+\nsubseteq\sigma\},\qquad
 K_-=\{\sigma\subseteq C:C_-\nsubseteq\sigma\}.
\]
For a simplex \(\sigma\) in a simplicial complex \(T\), define
\[
 \operatorname{lk}_T(\sigma)=\{\tau\in T:\tau\cap\sigma=\varnothing,\ \tau\cup\sigma\in T\}.
\]
For complexes on disjoint vertex sets, their join is \(K*L=\{\sigma\cup\tau:\sigma\in K,\ \tau\in L\}\). A bistellar flip replaces \(K_+*L\) by \(K_-*L\), where \(K_+\subseteq T\), the maximal simplices of \(K_+\) all have the same link \(L\) in \(T\), and the vertices of \(L\) are disjoint from \(C\). Explicitly, the replacement is
\[
 T'=\bigl(T\setminus((K_+*L)\setminus(K_-*L))\bigr)\cup(K_-*L).
\]
The move is allowed only when \(T'\) is a geometric triangulation belonging to the chosen vertex convention. The reverse replacement is also a flip.

For \(e\in\{\mathrm{all},\mathrm{sub}\}\), let \(G_e(A)\) have vertex set \(\mathcal T_e(A)\), with an edge between triangulations differing by one allowed bistellar flip. For each convention separately, the candidate question is whether
\[
 \forall d\in\{3,4\}\ \forall A\subset\R^d\text{ finite with }\operatorname{aff}(A)=\R^d,\qquad G_e(A)\text{ is connected}.
\]
Equivalently, every pair of triangulations in the chosen class must be joined by a finite sequence of allowed flips. The supplied record does not establish which candidate is intended, whether general position is required, or how nonspanning or infinite point sets should be treated.

The registry labels the problem open and supplies a dated assessment that connectivity fails in dimensions at least five while dimensions three and four remain open. These are supplied curation claims, not independently verified mathematical or literature conclusions.
\subsection{Short English statement}
Can any two triangulations of a point set in dimension \(3\) or \(4\) be connected by finitely many bistellar flips? The draft states explicit finite, full-dimensional versions, but the intended vertex convention and any general-position restriction require review. The supplied open status is unverified dated curation.
\subsection{Sources}
\begin{itemize}
\item Supplied local UnsolvedMath registry record, permanent ID \(30004519\), source code \texttt{catalogue entry}, in the Oberwolfach Reports Open Problems collection. Local registry references: \texttt{data/unsolvedmath/problems.json} and supplementary \texttt{catalog/problems.json}; no independent file inspection was performed. No individual proposer is supplied.
\item Supplied report citation: \emph{Discrete Geometry (2021)}, Oberwolfach Reports, DOI 10.4171/owr/2020/30, pp.\ 1469--1528. \url{https://doi.org/10.4171/owr/2020/30}. The registry gives proposed year 2020. The supplied citation is retained without independently resolving its year designation or locating the exact problem within the report.
\item Supplied literature source: Rote--Santos--Stump, \emph{Connectivity of Triangulation Flip Graphs}. \url{https://doi.org/10.1007/s00454-022-00436-2}.
\item Supplied literature source: TOPP, \emph{Flip Graph Connectivity in 3D}. \url{https://topp.openproblem.net/p28}.
\item Literature assessment dated exactly \texttt{2026-08-22}: the supplied curation states that bistellar flip-graph connectivity fails in dimensions at least five and remains open in dimensions three and four. This is supplied dated curation derived from the OWR record, not independently verified; its time-sensitive status and cited claims require independent checking.
\item The registry labels the statement \texttt{corrected\_verified}, but its verification note reports only repairs to formula delimiters, notation, grammar, and typography while preserving mathematical content and scope. This does not establish mathematical review of the present draft.
\end{itemize}
\end{document}
